\begin{description}
\item[(4.1)] To get the number of photons per second we have to multiply the
Flux by the area of the dish, to know the incoming energy per second, and
divide this by the energy of a photon.

For $\lambda_1 = 0.32\,mm = 3.2 \times 10^{-4}\,m$:

Energy of a photon:
\[ E = \frac{hc}{\lambda} = 6.2076 \times 10^{-22}\,Jules \]
% sic: the source spells "Jules" for Joules.
\hfill(1 point)

Area of the disk:
\[ A = \pi R^2 = 113\,m^2 \]

Number of photons per second:
\[ n_1 = \frac{flux \times A}{E} \approx 1820~photon/s \]
\hfill(1 point)

\item[(4.2)] Same calculation, just changing the energy of the individual
photons:

For $\lambda_2 = 8.6\,mm = 8.6 \times 10^{-3}\,m$

Energy of a photon:
\[ E = \frac{hc}{\lambda} = 2.31 \times 10^{-23}\,Jules \]
\hfill(1 point)

Number of photons per second:
\[ n_2 = \frac{flux \times A}{E} \approx 48900~photon/s \]
\hfill(1 point)

\item[(4.3)] Spatial resolution of a single telescope is given by:
\[ \theta = 1.22\,\frac{\lambda}{D} \]
where $D$ represents the diameter of the dish. This value will be given in
radians, so it must be converted to arcsec afterwards.

For a frequency of 74.9\,GHz, the corresponding wavelength is:
\[ \lambda = \frac{c}{f} = 4\,mm \]
\hfill(1 point)

And the spatial resolution:
\[ \theta = 1.22\,\frac{4\,mm}{12\,m} \approx 83.9~arcsec \]
\hfill(1 point)

\item[(4.4)] For an array the correct expression is:
\[ \theta = \frac{\lambda}{B} \quad\text{or}\quad
   \theta = \frac{\lambda}{2B} \]
\hfill(1 point)\\
being $B$ the longest baseline in the array.

So in this case:
\[ \theta = \frac{4\,mm}{16\,km} \approx 0.0516~arcsec
   \quad\text{or}\quad 0.0258~arcsec \]
\hfill(1 point)

\item[(4.5)] The SEFD is a characteristic flux of a system, found by
dividing the characteristic energy associated to the so-called temperature of
the antenna by its effective area:
\[ SEFD = \frac{2kT_{sys}}{A_e} \]
\hfill(1 point)

As no additional information is given about the effective area, the actual
physical area of the array should be used:
\[ A = 54\left(\pi \times 6^2\right) + 12\left(\pi \times 3.5^2\right)
   \approx 6569\,m^2 \]
\hfill(0.5 points)

Substituting the Boltzmann constant, the given temperatura of the antenna,
and converting the answer to Jansky we get:
% sic: "temperatura" appears in the source.
\[ SEFD \approx 290.5~Jy \]
\hfill(0.5 points)
\end{description}
